\documentclass{report}
\usepackage{amsmath,amssymb}
\setlength{\parindent}{0mm}
\setlength{\parskip}{1em}
\begin{document}
\begin{center}
\rule{6in}{1pt} \
{\large Eran Treister \\
{\bf Joint Full Waveform Inversion and Travel Time Tomography}}

Dept of Earth and Ocean Sciences \\ University of British Columbia \\ 2021 Main Mall \\ Vancouver \\ V6T-1Z4 \\ Canada
\\
{\tt erantreister@gmail.com}\\
Eldad Haber\end{center}

The full waveform inversion (FWI) in heterogeneous media is a highly
non-linear optimization problem. Using traditional methods like Gauss
Newton, the solution is highly sensitive to the initialization point. The
best results are often achieved
by a frequency continuation strategy where the process is initially
obtained with the low-frequency data only, and then gradually the high
frequency data
is added into the optimization. However, low frequency data is usually
missing, and in its absence the FWI process reaches a local minimum.
Travel time tomography, unlike FWI, is less sensitive to its
initialization and contains low
frequency features. In this work we propose a framework to jointly apply
FWI and travel time tomography, so that the resulting model is faithful
to both the
waveform and travel time data. Using this approach we also aim to relieve
the problem of missing low-frequency data. To this end, we minimize a sum
of the FWI and travel time tomography misfit functions, with
regularization.
Synthetic examples show that this strategy leads to a better recovery of
the underlying medium when low-frequency data is missing.


\end{document}
